\section{Repository structure and reproducibility}
\label{app:repo}

The codebase is organised as follows.

\begin{description}
\item[\code{morimura.py}] The inverter core: \code{Inverter} class with
softmax-parametric chain, analytic gradients (local + globals),
$\lambda$-cross-validated fit, and the synthetic $\S 6.1$ reproduction
driver.
\item[\code{congestion\_filter.py}] Early exploratory experiments:
snapshot-level
two-regime synthesis and detection.
\item[\code{per\_car\_detector.py}] Tier-2 per-trajectory log-likelihood-ratio detector.
\item[\code{sumo\_validation/}] The SUMO-side machinery:
\code{sumo\_to\_phase3.py} (network builder, trip reader, feature
extractor), \code{score\_seed\_big.py} and \code{contrast\_correlation.py}
(scorers, with the \code{--dataset \{sumo,xuancheng\}} switch),
the empirical and backoff $2$-step diagnostics
(\cref{sec:phase4-2step,sec:phase5-labour}), and
\code{fd\_check\_features.py} (the FD gradient gate). Exact script names
and log paths are indexed in \code{sumo\_validation/README.md}.
\item[\code{real\_data/}] The Xuancheng tier: the unmodified
\code{xuancheng.net.xml} from \cite{ma2026xuancheng}, the daily JSON
trip files, and \code{load\_xuancheng.py} which provides the loader,
the three regime-split helpers, and the OD-rebalancing extension.
\item[\code{papers/}] Source PDFs of all cited works.
\item[\code{research\_logbook.md}] The day-by-day project log,
with twenty experimental-cycle entries.
\end{description}

\section{Use of generative AI}\label{app:ai}

In accordance with the module's declaration requirement, generative AI
was used in this project as follows. LLM-based code assistants,
including Claude, Codex and Gemini, were used as pair-programming and
writing-support aids throughout the implementation and writeup phases:
drafting code (with the author's careful review of every line before
commit), writing unit-test scaffolding, deriving analytic gradient
formulas alongside the author, checking mathematical and experimental
arguments, and helping edit the report for clarity. The author retained
responsibility for the research design, interpretation, code accepted
into the repository, and final submitted text.

\section{Build and reproducibility}\label{app:build}

The project is reproducible from the repository via the README, with
the following minimum environment:

\begin{lstlisting}[language=bash]
# Python 3.14, see project README
python3 -m venv .venv
.venv/bin/python -m pip install numpy scipy matplotlib

# Tier 1: synthetic Morimura reproduction (~20 s)
.venv/bin/python morimura.py

# Tier 2: SUMO experiments (requires SUMO_HOME set)
.venv/bin/python sumo_validation/contrast_correlation.py \
    --seed 42 --features real --T 40 --maxiter 1500

# Tier 3: Xuancheng experiments (requires real_data/ populated)
.venv/bin/python sumo_validation/contrast_correlation.py \
    --dataset xuancheng \
    --day 2023-04-17,2023-04-18,2023-04-19,2023-04-20,2023-04-21 \
    --regime_split rush_offpeak --od_match 8 \
    --features real --T 10 --maxiter 1500

# Stronger-proxy Labour-Day stress test
DAYS=2023-04-28,2023-04-29,2023-04-30,\
2023-04-11,2023-04-12,2023-04-13,2023-04-14,\
2023-04-17,2023-04-18,2023-04-19,2023-04-20,2023-04-21
.venv/bin/python sumo_validation/contrast_correlation.py \
    --dataset xuancheng \
    --day "$DAYS" \
    --regime_split holiday_normal --od_match 64 \
    --features none --T 10 --maxiter 1500
\end{lstlisting}

\section{Metric and method primer}
\label{app:metrics-primer}

This appendix gives compact reminders for the mathematical objects,
optimisation routines and diagnostic metrics used in the main text. It
is not needed to follow the narrative, but it fixes the precise meaning
of terms that are easy to over-read from their names alone.

\subsection{Markov chain and transition kernel}
\label{app:markov-chain-primer}

A Markov chain is a stochastic process whose next state depends only on
its current state. In this dissertation the states are directed road
edges. If a vehicle is currently on edge $x$, the transition kernel
$\Pchain$ gives the probability of choosing each legal next edge:
\[
\pT(x'\mid x) = \Pr[X_{t+1}=x' \mid X_t=x].
\]
Each row of $\Pchain$ is therefore a probability distribution over the
legal successors of $x$: entries are non-negative and sum to one. The
``$1$-step'' assumption means that the model ignores the earlier path
history once $x_t$ is known. The empirical $2$-step diagnostic in
\cref{sec:phase4-2step} tests whether this simplification is too severe
by allowing the next edge to depend on $(x_{t-1},x_t)$.

\subsection{Stationary distribution, restart probability and hitting rate}
\label{app:stationary-primer}

The stationary distribution $\pi$ of a Markov chain is the long-run
fraction of time spent in each state if the chain is followed
indefinitely. It satisfies
\[
\pi^\top P = \pi^\top,\qquad \sum_x \pi(x)=1.
\]
The Morimura framework uses a restarted chain
\[
P = \beta \Pchain + (1-\beta)\bm{1}\pI^\top,
\]
where $\beta$ is the probability of continuing along the road-transition
kernel and $1-\beta$ is the probability of restarting from $\pI$. This
restart construction makes the chain well behaved and models finite trip
lengths.

A hitting rate $g(i,j)$ is a first-passage statistic: starting from
observed state $i$, how likely is the trajectory to hit observed state
$j$ later, discounted by the continuation probability $\beta$? In
principle hitting rates add path-order information beyond simple
station counts. In practice they are much noisier on sparse real-world
data because there are $|\Xo|^2$ station pairs to estimate.

\subsection{Inverse Markov-chain problem}
\label{app:inverse-mc-primer}

The forward Markov-chain problem asks: given $\Pchain$, what stationary
distribution and hitting rates does it imply? The inverse problem asks
the reverse: given partial observations of those statistics at a subset
$\Xo$, recover a plausible transition kernel. This is ill-posed without
extra structure, because many different kernels can produce similar
station counts at a few observation states. Morimura et al.'s method
\cite{morimura2013} supplies that structure by fitting a
softmax-parametric kernel and regularising the parameters.

\subsection{Softmax-parametric chain}
\label{app:softmax-primer}

A softmax turns arbitrary real-valued scores into probabilities. For a
road edge $x$ with legal successors $\mathrm{adj}(x)$, the model assigns
a score $s(x,x')$ to each possible next edge and sets
\[
\pT(x'\mid x)=
\frac{\exp(s(x,x'))}
{\sum_{y\in \mathrm{adj}(x)}\exp(s(x,y))}.
\]
This guarantees a valid probability distribution for every row of
$\Pchain$. It also lets the model share information through features:
the score can combine local edge-specific weights with global features
such as road class, speed, lane count or turn angle. In this thesis
``features=none'' means only local edge weights are used; ``features=real''
adds shared road-network features as additional parameters.

\subsection{Log-likelihood ratio score}
\label{app:llr-primer}

Given two fitted regime kernels $\hat{P}_A$ and $\hat{P}_B$, a trajectory
$\tau=(x_0,\ldots,x_T)$ is scored by the log-likelihood ratio
\[
\Lam(\tau)=
\sum_{t=0}^{T-1}\left[
\log \hat{p}_B(x_{t+1}\mid x_t)
-\log \hat{p}_A(x_{t+1}\mid x_t)
\right].
\]
Positive values mean the trajectory is more likely under regime $B$ than
regime $A$ according to the fitted chains; negative values mean the
opposite. The score is additive, so longer trajectories can accumulate
more evidence if the per-step differences are consistent.

\subsection{ROC curve and AUC}
\label{app:auc-primer}

The receiver-operating-characteristic (ROC) curve evaluates a scalar
score across all possible thresholds. For each threshold, it plots the
true-positive rate against the false-positive rate. The area under this
curve (AUC) is a threshold-free measure of separability:
\[
\AUC = \Pr[s(\tau_B) > s(\tau_A)]
\]
for a randomly chosen regime-$B$ trajectory and a randomly chosen
regime-$A$ trajectory, up to tie handling. An AUC of $0.5$ is chance;
$1.0$ is perfect ranking; values below $0.5$ indicate the score is
systematically reversed. AUC does \emph{not} prove correct chain
recovery: a classifier can rank regimes well because of aggregate bias,
OD mix or optimiser artefacts, which is why the thesis pairs AUC with
transition-contrast diagnostics.

\subsection{Pearson correlation}
\label{app:pearson-primer}

Pearson correlation measures linear alignment between two vectors after
subtracting their means:
\[
r(a,b)=
\frac{\sum_k (a_k-\bar a)(b_k-\bar b)}
{\sqrt{\sum_k(a_k-\bar a)^2}\sqrt{\sum_k(b_k-\bar b)^2}}.
\]
It lies in $[-1,1]$: $+1$ means perfect positive linear alignment, $0$
means no linear alignment, and $-1$ means perfect reversal. In this
thesis Pearson is computed on stacked legal outgoing-edge contrasts,
comparing
$\Delta_{\mathrm{fit}}=\hat{P}_B-\hat{P}_A$ with
$\Delta_{\mathrm{emp}}=P_{B,\mathrm{emp}}-P_{A,\mathrm{emp}}$. A positive
Pearson supports the claim that the fitted model moves probability mass
in the same directions as the full-observation empirical chain.

\subsection{Cosine similarity and collinearity}
\label{app:cosine-primer}

Cosine similarity measures directional alignment:
\[
\cos(a,b)=\frac{a^\top b}{\|a\|\,\|b\|}.
\]
It is insensitive to the vector magnitudes: two vectors pointing in the
same direction have cosine $+1$, orthogonal vectors have cosine $0$, and
opposite directions have cosine $-1$. We use per-row cosine at branching
states. For a road edge with multiple legal successors, the row vector
asks: did the fitted model increase and decrease probability on the same
outgoing turns as the empirical contrast? The median row cosine and the
fraction of rows with positive cosine are therefore local
chain-recovery diagnostics.

\subsection{OD matching / OD rebalancing}
\label{app:od-primer}

Origin-destination (OD) matching controls for where trips start and end.
Real-world regimes can differ because different people travel to
different places, not because comparable drivers choose different routes.
The Xuancheng loader clusters network nodes into $K$ spatial zones, bins
each trajectory by origin zone and destination zone, and then samples the
same number of trajectories from each regime inside every shared OD
cell. After this rebalancing, the two regime pools have the same coarse
OD distribution, so remaining AUC is closer to a within-OD route-choice
contrast. The sweep over $K$ checks whether this conclusion depends on
how coarse the OD zones are.

\subsection{Bootstrap confidence interval}
\label{app:bootstrap-primer}

A bootstrap estimates sampling uncertainty by repeatedly resampling the
observed test set with replacement. In the Xuancheng experiments the
bootstrap is stratified by regime: each replicate samples the same number
of regime-$A$ and regime-$B$ trajectories, recomputes the AUC, and stores
the result. The $2.5$th and $97.5$th percentiles form an empirical
$95\%$ confidence interval. These intervals capture test-trajectory
sampling noise conditional on the fitted chains, station set, OD match
and train/test split. They do not include uncertainty from refitting the
inverse-MC model or redrawing $\Xo$.

\subsection{L-BFGS-B optimiser, basins and multistart}
\label{app:lbfgs-primer}

L-BFGS-B is a quasi-Newton optimiser: it minimises a differentiable loss
using gradients and a limited-memory approximation to second-order
curvature. The ``B'' indicates support for simple bound constraints,
although the parameter fits here are effectively unconstrained because
the softmax maps real-valued scores to valid probabilities.

The inverse-MC loss is non-convex, so L-BFGS-B can converge to different
local optima, or basins, depending on the starting point. A multistart
check repeats the fit from several random initial parameters and selects
the model with the best validation log-likelihood. If multistart changes
the AUC or loss substantially, the single-start result may be an
optimiser artefact. If multistart fails to improve the detector, as in
the Labour-Day Xuancheng stress test, the weak fitted result is less
likely to be explained by a poor initialisation alone.

\subsection{Ridge regularisation and finite-difference gradient checks}
\label{app:regularisation-fd-primer}

Ridge, or Tikhonov, regularisation adds
$\frac{\lambda}{2}\|\theta\|^2$ to the loss. This discourages very large
parameter values and helps stabilise an under-constrained inverse
problem, especially when only a sparse subset of states is observed.

A finite-difference (FD) gradient check verifies an analytic gradient by
perturbing one parameter at a time:
\[
\frac{\partial L}{\partial \theta_k}
\approx
\frac{L(\theta+\epsilon e_k)-L(\theta-\epsilon e_k)}{2\epsilon}.
\]
The thesis uses FD checks as safety gates before trusting custom
gradient code, especially for the feature-augmented and sparse
hitting-rate variants.

\subsection{Gap closed}
\label{app:gap-closed-primer}

When the full-observation empirical reference has AUC above chance, the
gap-closed statistic reports how much of the available empirical signal
the fitted detector recovers:
\[
\mathrm{gap\ closed} =
\frac{\AUC_{\mathrm{fit}}-0.5}
{\AUC_{\mathrm{emp}}-0.5}.
\]
For example, if the empirical reference is $0.70$ and the fitted detector
is $0.60$, the fitted detector closes $(0.60-0.50)/(0.70-0.50)=50\%$ of
the above-chance gap. This statistic is undefined or uninformative when
the empirical reference is at chance.

\subsection{RMAE}
\label{app:rmae-primer}

Relative mean absolute error (RMAE) is used in the Tier~1 reproduction
because Morimura et al.'s original benchmark is a stationary-distribution
recovery task, not a two-regime detector. In this thesis it measures the
mean absolute error of the recovered stationary mass on unobserved states
relative to the scale of the true held-out counts. Lower is better. It is
therefore not directly comparable to AUC: RMAE evaluates reconstruction
accuracy, whereas AUC evaluates trajectory ranking between regimes.
